\setlength{\parindent}{2em}

\section{问题提出的背景}

\subsection{背景介绍}

\par 拖延症是困扰当代知识工作者的普遍性问题，尤其在大学生和年轻白领群体中表现突出。研究表明，约70\%的大学生存在不同程度的拖延行为，其中50\%以上的大学生认为拖延症对他们生活和态度造成了影响\cite{SJTY201919049}。拖延不仅导致任务完成质量下降、学习压力增大，更会引发焦虑、自责等负面情绪，形成“拖延-焦虑-更拖延”的恶性循环。从神经科学视角来看，拖延的本质是边缘系统对任务威胁的过度反应与前额叶皮层执行控制功能之间的失衡——当面对复杂、模糊或不熟悉的任务时，大脑将其识别为“威胁”并触发回避机制，而传统的任务管理工具恰恰未能针对这一心理机制进行有效干预\cite{denmark2010zeigarnik}。

\par 当前市场上的任务管理工具（如Todoist\cite{todoist}、滴答清单\cite{dida365}、番茄ToDo\cite{tomatodo}等）虽然功能丰富，但在解决拖延症核心痛点方面存在明显局限：（1）传统表单化界面增加了任务录入的认知负担，“开始五分钟最难”的问题未被有效解决；（2）机械式提醒对自律性较差的用户效果有限，缺乏情感化陪伴和个性化干预；（3）任务拆解依赖用户自身能力，对于“启动困难”型用户缺乏有效支持。

\par 随着大语言模型（LLM）技术的快速发展，生成式AI在任务规划、对话交互和个性化推荐方面展现出巨大潜力\cite{cao2023comprehensivesurveyaigeneratedcontent}。智能体框架也在协同演进，逐步从单一的对话模型演进为具备自主决策、工具调用和环境交互能力的复杂系统，AI技术正从“辅助生成”向“自主执行”的范式迁移。基于智能体的任务管理系统能够通过自然语言交互降低使用门槛，通过上下文感知提供精准干预，通过情感化陪伴增强用户粘性，为解决拖延症问题提供了新的技术路径。

\subsubsection{市场需求分析}

\par 效率管理工具市场具有较大的发展空间。在App Store上，具有数万下载量的任务管理软件有十余款，说明市场需求旺盛但竞争格局分散。目标用户群体主要为大学生和刚进入职场的年轻白领，这部分人群具有以下特征：（1）愿意为提升效率的新工具付费；（2）自我管理能力相对较弱，对AI辅助有较高接受度\cite{tresna2025fear}；（3）手机使用频率高，移动端交互是核心场景。

\par 从用户旅程角度分析，拖延行为的核心症结在于“第一步成本太高”\cite{klingsieck2013procrastination}。以大学生完成课程汇报为例：传统场景下，学生得知任务后往往因畏难情绪而推迟开始，直到截止日期临近才匆忙应对；使用传统日程工具虽能记录任务，但频繁的机械提醒反而加剧心理负担；而理想的AI助手应能在任务创建后主动提供拆解建议、渐进式知识输入和情感支持，降低启动门槛。

\subsubsection{技术发展趋势}

\par 随着人工智能技术从“被动响应”向“主动服务”的代际演进\cite{luo2025large}，对话式交互界面与智能体系统的发展为拖延症干预提供了全新的技术范式。传统的数字化干预工具往往依赖用户主动触发，存在交互割裂、反馈滞后等局限性；而基于大语言模型的智能体则能够通过自然语言交互构建拟人化的陪伴体验，实现全天候、沉浸式的行为引导，从而在心理机制层面有效降低任务启动阻力。

\par 在智能体系统构建层面，技术评估标准的确立标志着该领域正从概念探索走向工程化落地。Akshathala等人于2025年提出的智能体系统评估框架确立了四大核心支柱：大模型、会话记忆、工具调用和环境\cite{akshathala2025taskcompletionassessmentframework}。这一框架不仅为智能体任务管理系统的设计提供了理论支撑，更为拖延症干预系统的效能优化指明了技术路径——通过大模型的认知推理能力实现推理-执行循环\cite{yao2022react}，借助会话记忆维持干预连续性，利用工具调用拓展现实场景服务能力，并依托环境感知实现情境自适应交互。值得注意的是，随着Spring AI Alibaba\cite{spring_ai_alibaba}等国产开发框架的成熟，基于国产大模型\cite{liu2024deepseek}\cite{bai2023qwen}的智能体开发已具备完整的工具链支持，这不仅保障了用户数据的隐私安全，更为本土化行为干预模型的构建提供了技术底座。

\par 在人机协作范式创新方面，上下文感知计算与渐进式披露理论的融合应用，为低干扰设计提供了系统的指导原则。研究表明，移动端信息层级管理必须严格遵循认知负荷理论——在人类工作记忆有限容量的约束下，优先呈现最相关信息，将次要信息延迟展示或折叠处理\cite{sweller1988cognitive}。这种“按需供给”的交互逻辑与拖延症干预系统“低干扰、防挫败”的设计底线高度契合：当用户处于任务启动困难阶段时，系统应通过情境感知自动过滤非关键信息，仅呈现最小可行任务单元；随着用户认知资源的恢复，再逐步披露进阶支持功能。这种动态适配的信息架构不仅能有效避免认知超载引发的回避行为，更能通过渐进式任务呈现帮助用户建立“小步快跑”的成功体验，从而形成正向的行为强化循环。

\subsection{本研究的意义和目的}

\par 在数字化转型加速与知识经济蓬勃发展的时代背景下，知识工作者面临着前所未有的信息过载与多任务处理压力。拖延症作为一种普遍存在的心理与行为障碍，不仅表现为简单的“懒惰”，更深层次地折射出个体在面对复杂任务时的“启动困难”与执行过程中的“注意力断裂”等认知困境。传统的任务管理工具往往侧重于静态的清单罗列与时间提醒，缺乏对用户心理状态的感知与动态干预能力，难以从根本上解决拖延问题。因此，本研究旨在设计并实现一套基于智能体（Agent-based）的拖延症干预系统，通过融合人工智能、人机交互与行为心理学理论，探索一种全新的对话式任务管理模式，以期为知识工作者提供更具适应性与人性化的支持。

\subsubsection{理论价值：构建人机协同认知支持的理论框架}

\par 本研究致力于突破传统任务管理工具的静态交互局限，通过构建“用户意图-任务结构”动态映射模型，探索自然语言驱动的任务分解机制。具体而言，系统将引入意图识别算法与语义解析技术，将用户模糊的、非结构化的自然语言表达（如“我需要准备下周的报告”）自动转化为包含子任务、优先级与时间节点的结构化任务树。这一过程不仅涉及自然语言处理技术的优化，更触及人机协作式任务构建的核心逻辑——即智能体如何作为认知伙伴，辅助用户完成从抽象目标到具体行动的思维转化。此外，研究还将探索长期记忆持久化的系统化方案，通过构建用户行为画像与任务历史档案，使智能体能够基于用户习惯与偏好，提供个性化的任务建议与情境感知的干预策略。这些探索将为对话式交互设计理论贡献新的思路，推动生成式人工智能在认知辅助领域的理论深化，特别是在人机协同决策与动态任务规划方面提供新的研究视角。

\subsubsection{应用价值：打造融合对话交互与可视化管理的智能工具}

\par 在应用层面，本研究将结合智能体的记忆增强与工具调用能力，打破传统文本交互的局限。系统不仅能够与用户进行自然的对话，还能主动调用外部日历、待办清单或数据分析工具，并将对话流与可视化进度管理有效融合。例如，当用户通过对话表达“我想规划本周的工作”时，智能体不仅会生成任务列表，还会根据任务优先级自动生成甘特图或时间块视图，并实时更新任务进度。这种“对话即操作”的模式，旨在输出一种兼具交互自然性与执行推动力的解决方案。相比于现有的健康管理应用，本系统更加注重过程的引导与情感的共鸣，能够根据用户的实时反馈调整策略，如当检测到用户任务延迟时，自动触发鼓励性对话或提供时间管理建议。这些功能为数字健康管理领域提供了一套可落地、可复用的技术参考，具有广阔的市场应用前景，特别是在企业级知识管理与个人效能提升工具开发方面具有显著的应用潜力。

\subsubsection{社会价值：促进数字时代的个人福祉与社会效能}

\par 从更宏观的社会视角来看，拖延症不仅影响个人的学习与工作效率，还常常引发焦虑、内疚等负面情绪，进而对个体的心理健康与生活质量造成损害。本研究通过技术手段帮助拖延症用户建立健康的任务管理习惯，旨在降低因拖延带来的心理压力，提升个体的自我效能感与生活满意度。当广大知识工作者能够更高效、更从容地应对工作挑战时，整个社会的创新活力与生产效率也将得到相应提升。此外，本研究还关注数字技术对人类行为的积极引导作用，探索如何通过智能系统的设计，促进人与技术的和谐共生，避免技术依赖与数字倦怠。因此，本研究不仅具有技术先进性，更承载着促进数字时代个人福祉、构建健康数字工作环境的社会责任，为应对信息过载时代的普遍性挑战提供了有益的探索方向。

